\documentclass[11pt,a4paper]{article}
\usepackage{notes}

\title{Fourier Transforms in Geophysics II:\\
Spatial Sampling, Spectra}
\usepackage{fancyhdr}

\pagestyle{fancy}
\fancyhf{}
\lhead{Geophysics 3A: Potential Fields \& Geothermics}
\rhead{Week 7 Nevada Case Study}
\cfoot{\thepage}

\title{Week 7 --- Case Study: Estimating Basin Depth in Nevada using FFTs}
\author{Geophysics 3A: Potential Fields \& Geothermics}
\date{\today}

\begin{document}
\maketitle

\section{Purpose of the Demonstration}

This demonstration extends the basic ideas of Fourier transforms developed in Demo~1 to spatial potential-field data.  The focus is on understanding how sampling, filtering, and physical assumptions control what geological information can be recovered from gravity data.

The notebook progresses deliberately from:
\begin{itemize}
  \item synthetic geological models with known truth,
  \item to real Complete Bouguer gravity data from Nevada,
  \item to physically motivated removal of regional signals via isostatic correction,
  \item and finally to spectral interpretation and inversion.
\end{itemize}

\section{In-Class Demonstration Outline}

A suggested walkthrough is:

\begin{enumerate}
  \item \textbf{Nevada Bouguer gravity:} identify basin-scale stripes.
  \item \textbf{Isostatic correction:} remove long-wavelength crustal signal.
  \item \textbf{Spectral peaks:} identify dominant basin wavelengths.
  \item \textbf{Aliasing pipeline:} demonstrate failure of inversion under coarse sampling.
\end{enumerate}


\section{Spatial Fourier Transform}

For a uniformly sampled profile, the discrete Fourier transform gives
\begin{equation}
\hat{g}(k_n) = \sum_{j=0}^{N-1} g(x_j)\, e^{-i 2\pi k_n x_j},
\end{equation}
with discrete wavenumbers
\begin{equation}
k_n = \frac{n}{N\,\Delta x}.
\end{equation}

The amplitude spectrum $|\hat{g}(k)|$ shows how much power exists at each wavelength.
For potential fields, spectra typically show a \emph{red-noise} background:
power increases systematically toward longer wavelengths.

\section{Aliasing in the Spectral Domain}

When the sampling interval increases:
\begin{itemize}
  \item the Nyquist wavelength moves to longer scales,
  \item true basin-scale power migrates to incorrect longer wavelengths,
  \item spurious spectral peaks appear that look statistically significant.
\end{itemize}

The Fourier transform faithfully represents the sampled data,
even when that data are aliased.

\section{Airy Isostatic Gravity Correction}

The long-wavelength Bouguer gravity field over Nevada is dominated by
isostatic compensation of crustal thickness variations.
Under Airy isostasy, surface loads are compensated by a crustal root at depth $T_0$.

In the spectral domain, the gravity effect of a density interface at depth $T_0$ is
\begin{equation}
\hat{g}_{\text{root}}(\mathbf{k})
=
2\pi G\, \Delta\rho \, \hat{h}(\mathbf{k}) \, e^{-|\mathbf{k}| T_0},
\end{equation}
where:
\begin{itemize}
  \item $\Delta\rho$ is the density contrast across the Moho,
  \item $e^{-|\mathbf{k}|T_0}$ is the upward-continuation filter.
\end{itemize}

Subtracting this predicted root contribution from the observed Bouguer gravity yields
the \emph{isostatic residual}:
\begin{equation}
g_{\text{iso}} = g_{\text{Bouguer}} - g_{\text{root}}.
\end{equation}

This is a physically motivated, band-limited removal of long-wavelength power,
not a statistical detrend.

\section{Spectral Interpretation of Nevada Gravity}

After isostatic correction:
\begin{itemize}
  \item the mean gravity is near zero,
  \item long-wavelength crustal signals are suppressed,
  \item basin-scale anomalies (20--60~km) are enhanced.
\end{itemize}

Mean E--W and N--S power spectra reveal:
\begin{itemize}
  \item elevated E--W power at 20--40~km wavelengths,
  \item longer characteristic wavelengths in the N--S direction,
  \item a red-noise background consistent with potential-field theory.
\end{itemize}

\section{Statistical Significance of Spectral Peaks}

Because adjacent profiles are spatially correlated,
the number of independent spectral estimates is smaller than the number of rows or columns.
An effective number of degrees of freedom is estimated from the spatial autocorrelation length.

After fitting and removing a power-law background,
a chi-squared test is used to identify statistically significant spectral peaks.
For Nevada:
\begin{itemize}
  \item the dominant E--W peak lies near $\lambda \approx 20$--25~km,
  \item consistent with Basin-and-Range basin spacing.
\end{itemize}

\section{High-Pass Filtering and Inversion}

Before inverting gravity for sediment thickness,
residual long-wavelength power must be removed.
A raised-cosine high-pass filter is applied:
\begin{itemize}
  \item wavelengths $> 100$~km are blocked,
  \item wavelengths $< 60$~km are passed,
  \item a smooth taper prevents Gibbs ringing.
\end{itemize}

The filtered gravity is inverted using the Parker--Oldenburg method,
which relates gravity to interface relief in the spectral domain.

\section{Impact of Aliasing on Inversion}

Repeating the full processing pipeline on decimated data shows that:
\begin{itemize}
  \item when the Nyquist wavelength exceeds basin width,
  \item sediment thickness maps degrade rapidly,
  \item artefacts appear that have no geological meaning.
\end{itemize}

No inversion algorithm can recover information lost to aliasing.

\section{Sediment Thickness Inversion and Spectral Filtering}

After removal of the long-wavelength isostatic signal, the remaining gravity field
primarily reflects density variations in the shallow crust.
In the Basin-and-Range Province, these variations are dominated by
low-density sediment filling extensional basins over denser crystalline basement.
We now describe how this signal is converted into a sediment thickness model,
and why additional spectral filtering is required before inversion.

\subsection{Physical Meaning of the Sediment Thickness Model}

The sediment thickness model $h(x,y)$ represents the vertical thickness of
low-density basin fill that replaces higher-density basement rock.
The relevant density contrast is
\begin{equation}
\Delta\rho = \rho_{\text{basement}} - \rho_{\text{sediment}},
\end{equation}
which is taken here to be approximately $670~\text{kg m}^{-3}$.

Under this convention:
\begin{itemize}
  \item negative gravity anomalies correspond to thick sedimentary basins,
  \item positive anomalies correspond to basement highs or exposed ranges.
\end{itemize}

The model treats sediment as a single layer with uniform density overlying basement.
This is a deliberate simplification.
Gravity is most sensitive to the largest density contrast,
and cannot resolve fine-scale stratigraphy within the sediment column.
The resulting sediment thickness therefore captures the dominant mass deficit
associated with basin fill, rather than detailed internal layering.

\subsection{Forward Gravity Model for an Interface}

For an interface with mean depth $z_0$, relief $h(x,y)$, and density contrast
$\Delta\rho$, the linearized gravity response in the Fourier domain is
\begin{equation}
\hat{g}(\mathbf{k})
=
2\pi G\,\Delta\rho \;
e^{-|\mathbf{k}| z_0}
\;
\hat{h}(\mathbf{k}),
\end{equation}
where $\mathbf{k}$ is the horizontal wavenumber vector and
$e^{-|\mathbf{k}| z_0}$ is the upward-continuation kernel.

This expression highlights two fundamental properties of gravity data:
\begin{itemize}
  \item short wavelengths decay rapidly with depth,
  \item gravity observations are inherently band-limited by source depth.
\end{itemize}

\subsection{Instability of Direct Inversion}

Formally inverting the linear relation gives
\begin{equation}
\hat{h}(\mathbf{k})
=
\frac{\hat{g}(\mathbf{k})}
{2\pi G\,\Delta\rho}
\;
e^{|\mathbf{k}| z_0}.
\end{equation}

The exponential factor $e^{|\mathbf{k}| z_0}$ amplifies high-wavenumber components.
As a result:
\begin{itemize}
  \item noise,
  \item discretization error,
  \item aliasing from inadequate sampling,
\end{itemize}
are dramatically magnified.
This makes naive downward continuation unstable and physically meaningless.

\subsection{The Parker--Oldenburg Method}

The Parker--Oldenburg method (Parker, 1973; Oldenburg, 1974) addresses this instability
by treating the gravity--interface relationship as weakly nonlinear
and solving iteratively in the spectral domain.

In practice, the method:
\begin{enumerate}
  \item starts with an initial guess for the interface relief $h(x,y)$,
  \item forward-calculates the gravity anomaly,
  \item compares predicted and observed gravity,
  \item updates the interface estimate,
  \item iterates until convergence.
\end{enumerate}

Each iteration uses fast Fourier transforms, with computational cost
$O(N\log N)$, making the method efficient even for large grids.
The result is a stable estimate of interface relief, provided the input gravity
is appropriately filtered.

\subsection{Why High-Pass Filtering Is Required}

Even after isostatic correction, the gravity field may still contain
long-wavelength signals arising from:
\begin{itemize}
  \item imperfect isostatic compensation,
  \item lithospheric flexure,
  \item deeper crustal or mantle density variations.
\end{itemize}

If these wavelengths are left in the data,
the inversion will interpret them as extremely thick sediment layers,
producing unphysical basin depths.

High-pass filtering therefore encodes a geological prior:
\begin{quote}
sediment thickness varies on wavelengths comparable to basin widths,
not on lithospheric scales.
\end{quote}

\subsection{Raised-Cosine High-Pass Filter}

A sharp spectral cutoff produces Gibbs ringing in the spatial domain,
manifesting as oscillations and spurious basin edges.
To avoid this, a smooth raised-cosine (Hann) filter is applied:
\begin{equation}
H(k) =
\begin{cases}
0, & k \le k_{\text{lo}}, \\
\frac{1}{2}
\left[
1 - \cos\!\left(
\pi \frac{k - k_{\text{lo}}}{k_{\text{hi}} - k_{\text{lo}}}
\right)
\right],
& k_{\text{lo}} < k < k_{\text{hi}}, \\
1, & k \ge k_{\text{hi}}.
\end{cases}
\end{equation}

In this application:
\begin{itemize}
  \item wavelengths $> 100~\text{km}$ are fully suppressed,
  \item wavelengths $< 60~\text{km}$ are fully passed,
  \item intermediate wavelengths are smoothly tapered.
\end{itemize}

This filter removes residual regional power while preserving basin-scale structure,
and ensures that the Parker--Oldenburg inversion remains stable and geologically
plausible.

\subsection{Interpretive Context}

It is important to emphasize that:
\begin{itemize}
  \item the inversion is not unique,
  \item the sediment thickness model is conditional on sampling, filtering,
        and density assumptions,
  \item no inversion can recover wavelengths shorter than the Nyquist limit.
\end{itemize}

The success of the Parker--Oldenburg inversion in this case reflects the fact that
sampling, isostatic correction, and spectral filtering have already constrained
the problem to a physically interpretable band of wavelengths.

\subsection{Key Takeaway}

The Parker--Oldenburg method is best understood as a physically motivated,
regularized downward continuation.
It produces meaningful sediment thickness estimates only when the gravity data
have been appropriately sampled, filtered, and interpreted.

\end{document}
